\chapter{Prototype Architecture}

To validate the feasibility of the Behavioral Engine for Demand (BED) in a
real-world setting, a lightweight, privacy-safe prototype was developed. The
prototype demonstrates how BED can be deployed using modern serverless
infrastructure, enabling real-time event logging, treatment assignment, and
behavioral state updates without requiring integration with point-of-sale (POS)
systems or sensitive user data. This chapter describes the architecture,
components, and data flow of the BED v3 prototype.

\section{Design Principles}

The prototype was designed according to four principles:

\begin{enumerate}
    \item \textbf{Safety}: No personally identifiable information (PII) is
    collected or stored.
    \item \textbf{Simplicity}: The system uses minimal infrastructure to reduce
    operational overhead.
    \item \textbf{Scalability}: Serverless components automatically scale with
    traffic.
    \item \textbf{Modularity}: Each component can be replaced or extended
    independently.
\end{enumerate}

These principles ensure that the prototype is suitable for academic research,
field experiments, and early-stage deployments.

\section{System Overview}

The prototype consists of four components:

\begin{enumerate}
    \item \textbf{QR Code Trigger}: Users scan a QR code that encodes a URL.
    \item \textbf{Static Web Page}: A lightweight HTML page executes a POST
    request to the backend.
    \item \textbf{Serverless Backend}: A cloud function assigns treatments and
    logs events.
    \item \textbf{Firestore Event Log}: A structured database stores synthetic
    event records.
\end{enumerate}

This architecture mirrors the closed-loop structure of BED while remaining
simple enough for rapid experimentation.

\section{QR Code Trigger}

The prototype begins with a QR code printed on a physical card or displayed on a
screen. The QR code encodes a URL pointing to the static web page. When scanned,
the user’s device opens the page in a browser. No user identifiers are
collected; each scan is treated as an independent event.

\section{Static Web Page}

The static page is implemented as a minimal HTML file. Upon loading, it
automatically sends a POST request to the serverless backend:

\begin{itemize}
    \item event type (e.g., ``scan''),
    \item timestamp,
    \item optional synthetic identifiers (if enabled).
\end{itemize}

The page itself contains no interactive elements and performs no tracking. Its
sole purpose is to trigger the backend event pipeline.

\section{Serverless Backend}

The backend is implemented as a cloud function that performs three tasks:

\begin{enumerate}
    \item \textbf{Treatment Assignment}: Randomly assigns users to control or
    treatment groups.
    \item \textbf{Incentive Allocation}: Determines whether a BED token or
    synthetic reward is issued.
    \item \textbf{Event Logging}: Writes a structured event record to Firestore.
\end{enumerate}

The backend does not store IP addresses, device identifiers, or cookies. All
events are synthetic and anonymized.

\section{Event Logging Schema}

Events are stored in a Firestore collection using the following schema:

\begin{itemize}
    \item \texttt{timestamp}: time of event,
    \item \texttt{event\_type}: type of action (e.g., scan),
    \item \texttt{treatment\_group}: control or treatment,
    \item \texttt{bed\_tokens}: number of tokens assigned,
    \item \texttt{node\_id}: synthetic state identifier,
    \item \texttt{action\_id}: synthetic action identifier.
\end{itemize}

This schema mirrors the event log structure used in the simulation environment,
enabling direct comparison between simulated and prototype data.

\section{Closed-Loop Alignment with BED}

Although simplified, the prototype preserves the essential closed-loop structure
of BED:

\begin{enumerate}
    \item An event is triggered (QR scan).
    \item The backend assigns a treatment.
    \item The event is logged.
    \item The log can be used to update behavioral state offline.
\end{enumerate}

This demonstrates how BED can be deployed in real-world environments without
requiring invasive data collection or complex integrations.

\section{Scalability and Deployment}

The prototype is fully serverless:

\begin{itemize}
    \item The static page is hosted on a CDN.
    \item The backend scales automatically with traffic.
    \item Firestore handles concurrent writes without manual provisioning.
\end{itemize}

This architecture supports thousands of events per second with minimal cost,
making it suitable for field experiments and pilot deployments.

\section{Limitations}

The prototype is intentionally minimal and has several limitations:

\begin{itemize}
    \item It does not track user identity or longitudinal behavior.
    \item It does not integrate with POS or transactional systems.
    \item It does not perform real-time state updates.
\end{itemize}

These limitations are acceptable for early-stage experimentation and can be
addressed in future iterations.

\section{Summary}

The BED v3 prototype demonstrates a practical, safe, and scalable approach to
deploying closed-loop incentive systems. It validates the feasibility of the
BED architecture and provides a foundation for real-world experimentation. The
next chapter describes the experimental design used to evaluate BED policies.
